\chapter{1977:Ilya Prigogine}

\label{chap:chemistry1977}

The Nobel Prize in Chemistry for the year 1977 was awarded to Ilya Prigogine.

\textbf{Motivation:}

for his contributions to non-equilibrium thermodynamics, particularly the theory of dissipative structures

\section{Ilya Prigogine}

\subsection*{Early Life and Education}

Ilya Prigogine was born on 1917-01-25 in Moscow, Russia.

\subsection*{Career and Major Research}

Prigogine studied at the Free University of Brussels, where he received his doctorate and later became a professor. He also held a professorship at the University of Texas at Austin. He developed the thermodynamics of irreversible processes and introduced the concept of dissipative structures, ordered states that form in systems far from equilibrium.

\subsection*{Nobel Prize-winning Work}

The work for which Ilya Prigogine was awarded the Nobel Prize in Chemistry in 1977 was described by the Nobel Committee as: "for his contributions to non-equilibrium thermodynamics, particularly the theory of dissipative structures".

Prigogine showed that open systems exchanging energy and matter with their surroundings can spontaneously organize into stable, ordered structures such as those seen in oscillating chemical reactions.

\subsection*{Significance and Impact}

Dissipative structures connected chemistry to biology by explaining how order can arise in non-equilibrium systems, influencing fields from chemical kinetics to ecology.

\subsection*{Later Career and Honors}

Prigogine worked in both Brussels and Austin for the rest of his career. He died on 2003-05-28 in Brussels, Belgium.

\subsection*{Legacy}

Prigogine's ideas on self-organization in far-from-equilibrium systems remain influential across the natural and social sciences.

\subsection*{Modern Developments}

Prigogine's theory of dissipative structures explained how order can emerge from chaos in far-from-equilibrium systems, providing a scientific foundation for understanding self-organization in chemistry and biology. His ideas are applied in pattern formation, chemical oscillations such as the Belousov–Zhabotinsky reaction, and the thermodynamics of life itself.

\subsection*{Selected Publications}

"Introduction to Thermodynamics of Irreversible Processes" (1955)

"From Being to Becoming: Time and Complexity in the Physical Sciences" (1980)

\clearpage

\section*{Chapter Summary}

The 1977 prize recognized Ilya Prigogine for his contributions to non-equilibrium thermodynamics. His theory of dissipative structures explains how order can emerge in systems far from equilibrium.

